\chapter{اللقاء الثاني عشر: بداية المجلد الثاني - قصة استسقاء موسى وما بعدها}

\section{تأويل (وإذ استسقى موسى لقومه)}
السلام عليكم ورحمة الله وبركاته. الحمد لله رب العالمين، والصلاة والسلام على نبينا محمد وعلى آله وصحبه أجمعين.
وصلنا بحمد الله تبارك وتعالى إلى بداية المجلد الثاني، من أول قول الله تبارك وتعالى: \begin{ayah}وَإِذِ اسْتَسْقَىٰ مُوسَىٰ لِقَوْمِهِ فَقُلْنَا اضْرِب بِّعَصَاكَ الْحَجَرَ ۖ فَانفَجَرَتْ مِنْهُ اثْنَتَا عَشْرَةَ عَيْنًا ۖ قَدْ عَلِمَ كُلُّ أُنَاسٍ مَّشْرَبَهُمْ\end{ayah}.

\isnad{قال الإمام الطبري رحمه الله:} «يعني جل ثناؤه بقوله \textit{وإذ استسقى موسى}: وإذ استسقانا موسى لقومه، أي سألنا أن نسقي قومه ماء. فترك ذكر المسؤول ذلك، والمعنى الذي سأل موسى، إذ كان فيما ذكر من الكلام الظاهر دلالة على معنى ما ترك وحذف... وكذلك قوله \textit{قد علم كل أناس مشربهم} إنما معناه: قد علم كل أناس منهم مشربهم. فترك ذكر (منهم) لدلالة الكلام عليه».

ثم ساق \rawi{الطبري} الأخبار في تفصيل ذلك، \isnad{فروى بإسناده عن} \rawi{قتادة} و\rawi{ابن عباس} و\rawi{مجاهد} و\rawi{ابن زيد} وغيرهم أن ذلك كان في التيه، حين اشتكوا إلى نبيهم الظمأ، وأن الحجر كان حجراً مربعاً يُحمل معهم، تنفجر منه اثنتا عشرة عيناً لكل سبط عين.

\begin{fawaid}[عدم توقف العبرة على تفاصيل الإسرائيليات]
نلاحظ أن هذه الآثار المروية المذكورة فيها تفاصيل كنوع الحجر (مربع، أو مثل رأس الشاة) وكيفية الانفجار. وواضح جداً أن هذه الآثار مأخوذة عن بني إسرائيل. وكما بيّنا سابقاً، فإن ما جاء في كتاب الله هو ما تؤخذ منه العبرة، وهذه التفاصيل الزائدة لم تُذكر في كتاب الله ولا في حديث رسول الله ﷺ، فلا يُبنى عليها حكم. العبرة في القصة تتم حتى مع عدم العلم بهذه التفاصيل؛ فتروى ولا تصدق ولا تكذب.
\end{fawaid}

\separator

\section{تأويل (كلوا واشربوا من رزق الله ولا تعثوا في الأرض مفسدين)}
\isnad{قال أبو جعفر:} «أخبر جل ثناؤه أنه أمرهم بأكل ما رزقهم في التيه من المن والسلوى، وبشرب ما فجر لهم من الماء من الحجر المتعاور الذي لا قرار له في الأرض... ثم تقدم جل ثناؤه إليهم مع إباحته لهم ما أباح وإنعامه عليهم... بالنهي عن السعي في الأرض فساداً، والعثا فيها استكباراً. فقال تعالى ذكره لهم: \textit{ولا تعثوا في الأرض مفسدين}. أي لا تطغوا، ولا تسعوا في الأرض مفسدين».

وأصل (العَثَا): شدة الإفساد، بل هو أشد الإفساد، يقال منه: عثي فلان في الأرض، إذا تجاوز في الإفساد إلى غايته.

\begin{fawaid}[استعمال النعم في معصية الله]
إذا أنعم الله تبارك وتعالى على عبده بنعمة، فإن أشد الجحود أن تُستعمل نعمة الله في معصية الله. ولذلك كان توجيه الخطاب لهم بعد تذكيرهم بالمن والسلوى وانفجار الماء أن لا يفسدوا في الأرض. فالشكر يقتضي استعمال النعمة في مرضاة المنعِم.
\end{fawaid}

\begin{fawaid}[الآيات والنعم لا توجب الهداية بنفسها]
من المعاني التي يكررها هذا الموضع أن كثرة الآيات والنعم لا تكفي وحدها في حصول الإيمان والشكر؛ فقد رأى بنو إسرائيل من آيات الله وإنعامه ما لم يره كثير من الناس، ثم قابل فريق منهم ذلك بالكفر والعصيان والاعتراض. وفي هذا تسلية للداعية، وتعليم له أن الهداية بيد الله، وأن وظيفته البلاغ والبيان، لا حمل القلوب على الإيمان قهراً.
\end{fawaid}

\separator

\section{تأويل (وإذ قلتم يا موسى لن نصبر على طعام واحد)}
\isnad{قال أبو جعفر رحمه الله:} «قد دللنا فيما مضى على معنى الصبر وأنه كف النفس وحبسها عن الشيء. فمعنى الآية إذاً: واذكروا إذ قلتم يا معشر بني إسرائيل لن نطيق حبس أنفسنا على طعام واحد. وذلك الطعام الواحد هو ما أخبر الله جل ثناؤه أنه أطعمهموه في تيههم، وهو السلوى... فاسأل لنا ربك يخرج لنا مما تنبت الأرض...».

ثم ناقش \rawi{الطبري} الخلاف النحوي في قوله: \textit{يخرج لنا مما تنبت الأرض}، وردّ على من زعم (كالأخفش) أن (مِن) هنا زائدة ملغاة بمعنى (ما تنبت الأرض)، ورجح أنها للتبعيض، أي: يخرج لنا بعض ما تنبت الأرض.

\subsection{القول في الفوم والترجيح في تأخير الأسانيد}
اختلفوا في معنى (الفُوم)؛ فقال بعضهم هو الحنطة والخبز، وهو مروي عن \rawi{عطاء} و\rawi{مجاهد} و\rawi{قتادة} و\rawi{ابن عباس} و\rawi{ابن زيد}. وقال آخرون: هو الثوم.

\begin{fawaid}[تقديم الأسانيد وتأخيرها عند الطبري]
من الأبحاث المهمة في منهج \rawi{الطبري} ملاحظة أنه أحياناً يؤخر ذكر الأثر عن الصحابي البارز كـ \rawi{ابن عباس} بعد إيراده لأقوال التابعين. والسبب الغالب في ذلك هو نظره في صحة الإسناد؛ فإذا كان الإسناد إلى الصحابي فيه شيء من الضعف أو الانقطاع أخّره، وقدّم القول المروي بأصح إسناد ولو كان إلى من هو دونه في الطبقة.
\end{fawaid}

\separator

\section{تأويل (اهبطوا مصراً) ومنهج الطبري عند الاختلاف}
\isnad{قال أبو جعفر:} «ثم اختلفت القَرَأَة في قراءة قوله: \textit{اهْبِطُوا مِصْرًا}. فقرأه عامة القرأة: اهبطوا مِصْراً بتنوين (مصر) وإجرائه. وقرأه بعضهم بترك التنوين وحذف الألف منه».

ثم بيّن \rawi{الطبري} أن من نوّنه قصد مِصْراً من الأمصار، ومن لم ينونه قصد به مِصْر المعروفة التي خرجوا منها (مصر فرعون). واختلف المفسرون تبعاً لذلك هل المقصود الشام أم مصر فرعون؟ وذكر حجة كل فريق.

\begin{fawaid}[منهج الطبري في عدم الجزم في موضع الظن]
بعد استعراض الخلاف، وضع \rawi{الطبري} قاعدة في التعامل مع الاحتمالات المتكافئة قائلاً: «والذي نقول به في ذلك أنه لا دلالة في كتاب الله جل ثناؤه على الصواب من هذين التأويلين، ولا خبر به عن الرسول ﷺ يقطع مجيئه العذر، وأهل التأويل فيه متنازعون... فأولى الأقوال في ذلك عندنا بالصواب أن يقال: أمرهم أن يهبطوا قراراً من الأرض تنبت ما سألوا... وجائز أن يكون ذلك مصر، وجائز أن يكون الشام». وهذا من أحسن ما يفعله المحقق؛ ألا يجزم في موضع ظن تتكافأ فيه الأدلة، بل يعمم ما عممه القرآن.
\end{fawaid}

\separator

\section{تأويل (وضربت عليهم الذلة والمسكنة وباءوا بغضب من الله)}
\isnad{قال أبو جعفر:} «يعني جل ثناؤه بقوله \textit{وضربت}: أي فرضت ووضعت عليهم الذلة وألزموها... والذلة هي الصغار الذي أمر الله عباده المؤمنين ألا يعطوهم أماناً على القرار على ما هم عليه... حتى يعطوا الجزية عن يد وهم صاغرون. وأما المسكنة: مسكنة الفاقة والحاجة، وهي خشوعها وذلها... وأما \textit{باءوا بغضب}: انصرفوا ورجعوا متحملين غضباً من الله».

\begin{fawaid}[إثبات أفعال الله والرد على الجهمية]
تطرق \rawi{الطبري} في هذا الموضع وفي نظائره إلى إثبات أفعال الله تبارك وتعالى كالغضب والاستهزاء والسخرية والمكر، وأثبتها لله كما أخبر تبارك وتعالى عن نفسه، معتنياً في كتابه بالرد على المعتزلة والجهمية الذين عطّلوا معاني أفعال الله وحرّفوها بادعاء التنزيه.
\end{fawaid}

\separator

\section{تأويل (ويقتلون النبيين بغير الحق) وأصل تسمية "النبي"}
\isnad{قال أبو جعفر:} «يعني بذلك فعلنا بهم ذلك من أجل أنهم كانوا يجحدون حجج الله... ويقتلون رسل الله الذين بعثهم لإنباء ما أرسلهم به عنه... وهم جِماع واحدهم (نبي) بغير همز، وأصله الهمز، لأنه من أَنْبَأَ عن الله، فهو يُنبِئ عنه إنباء، وإنما الاسم منه مُنْبِئ، ولكنه صُرف وهو مُفعِل إلى فَعيل...».
ثم ذكر قولاً آخر بأن النبي والنبوة مأخوذان من النَّبْوة (وهي المكان المرتفع)، وكأن أصل النبي الطريق الظاهر المستبين.

\begin{fawaid}[أصل تسمية النبي لغة]
الأقرب والأقوى هو القول الأول؛ أن "النبي" مأخوذ من "النبأ"، لأنه مُنْبَأٌ من الله (يوحى إليه) ومُنْبِئٌ عن الله (مخبر). وخاصة النبي أنه يُنبئه العليم الخبير، وهذا يتضح في قوله تعالى: \textit{نبأني العليم الخبير}.
\end{fawaid}

\begin{fawaid}[دلالة القتل بغير حق]
قوله تعالى: \textit{بغير الحق} ليس معناه استشكال أن هناك قتلاً للأنبياء "بحق"! بل المراد وصف فعلهم بأنه مجرد عن أي إذن أو حجة، فقد جمعوا بين الكفر بالأنبياء والاعتداء غير المأذون فيه عليهم بقتلهم ظلماً.
\end{fawaid}

\separator

\section{تأويل (ذلك بما عصوا وكانوا يعتدون) والعبرة من قصص القرآن}
\isnad{قال أبو جعفر:} «المعنى: ذلك بعصيانهم وكونهم معتدين، والاعتداء تجاوز الحد... فمعنى الكلام فعلت بهم ما فعلت من ذلك بما عصوا أمري وتجاوزوا حدي».

\begin{fawaid}[أسباب نزول العقوبات والعبرة من القصص]
من الملاحظات الدقيقة في كتاب الله أن الله إذا ذكر عقاباً للكفار أو الظالمين في الدنيا أو الآخرة أرجع السبب إليهم \textit{بما عصوا}، \textit{فبظلم من الذين هادوا}، \textit{ذلك جزيناهم ببغيهم}؛ فالله لا يظلم مثقال ذرة. 
كما أن القصة سيقت هنا موعظةً وتعليماً للمؤمنين بألا يسلكوا مسلك اتباع الأنبياء الذين كفروا بعد إذ رأوا الآيات، وعزاءً لرسول الله ﷺ بأن جحود معاصريه من اليهود ليس بدعاً، بل هو امتداد لسلفهم الذين قتلوا الأنبياء.
\end{fawaid}

\separator

\section{تأويل (إن الذين آمنوا والذين هادوا والنصارى والصابئين...)}
\isnad{قال أبو جعفر:} «أما الذين آمنوا فهم المصدقون رسول الله ﷺ... وأما الذين هادوا فهم اليهود، ومعنى هادوا تابوا... والنصارى جمع واحدهم نصران كنشوان وسكران... وسموا نصارى لنصرة بعضهم بعضاً... والصابئون جمع صابئ وهو المُسْتَحْدِثُ سوى دينه ديناً، كالمرتد من أهل الإسلام...».

وذكر اختلافهم في من يلزمه اسم الصابئين؛ فقيل هم قوم لا دين لهم (بين اليهود والمجوس)، وقيل قوم يعبدون الملائكة ويصلون القبلة، وقيل طائفة من أهل الكتاب. 
ثم قال في تأويل \begin{ayah}مَنْ آمَنَ بِاللَّهِ وَالْيَوْمِ الْآخِرِ وَعَمِلَ صَالِحًا\end{ayah}: «من آمن منهم بمحمد ﷺ وبما جاء به واليوم الآخر، ويعمل صالحاً فلم يبدل ولم يغير حتى توفي على ذلك...».

\begin{fawaid}[الإيمان قول وعمل ورد المرجئة]
فسّر \rawi{الطبري} الإيمان هنا بالتصديق والثبات والعمل الصالح وعدم التبديل حتى الوفاة. وهذا يرد على من ظن أن الإيمان مجرد تصديق، أو أن قوله تعالى \textit{آمنوا} يغني عن الثبات والعمل. فالإيمان شُعب ودرجات، ولا يكون مقبولاً بمجرد التصديق دون اتباع.
\end{fawaid}

\subsection{رواية سلمان الفارسي وضعف بعض الأسانيد للغرابة}
ساق \rawi{الطبري} في سبب نزول هذه الآية القصة الطويلة المشهورة لرحلة \rawi{سلمان الفارسي} رضي الله عنه في البحث عن الحق، وصحبته للرهبان، وكيف سأل النبي ﷺ عنهم، فأنزل الله هذه الآية.

\begin{fawaid}[رد الإسناد لغرابة المتن ونكارته]
في تعليقه على بعض الآثار (كرواية تسمية النصارى نسبة لقرية الناصرة)، قال الطبري: «وقد ذُكر عن ابن عباس من طريق غير مُرْتَضَى...». الطبري قد يضعف الإسناد إذا وجد أن المتن منكر أو غريب وليس له شاهد، حتى وإن مرّر نفس الإسناد في مواضع أخرى لا استنكار في متنها.
\end{fawaid}

\separator

\section{تأويل (وإذ أخذنا ميثاقكم ورفعنا فوقكم الطور خُذوا ما آتيناكم بقوة)}
\isnad{قال أبو جعفر:} «الميثاق المِفْعال من الوثيقة... وقوله \textit{ورفعنا فوقكم الطور}، أما الطور فإنه الجبل في كلام العرب، وقيل إنه جبل بعينه... وتأويل الكلام: خذوا ما افترضنا عليكم في كتابنا من الفرائض، فاقبلوه واعملوا باجتهاد منكم في أدائه من غير تقصير ولا توانٍ، وذلك هو معنى أخذه إياه بقوة وبجد».

\isnad{وقال في تأويل} \begin{ayah}ثُمَّ تَوَلَّيْتُم مِّن بَعْدِ ذَٰلِكَ\end{ayah}: «أي خالفوا ما كانوا وعدوا الله... ونبذوا ذلك وراء ظهورهم. \begin{ayah}فَلَوْلَا فَضْلُ اللَّهِ عَلَيْكُمْ وَرَحْمَتُهُ لَكُنتُم مِّنَ الْخَاسِرِينَ\end{ayah} أي فلولا أن الله تفضل عليكم بالتوبة بعد نكثكم الميثاق، وأنعم عليكم بالإسلام ورحمته التي رحمكم بها... لكنتم الباخسين أنفسهم حظوظها».

بهذا نختم درسنا اليوم، ونكمل غداً بإذن الله. والسلام عليكم ورحمة الله وبركاته.
